\textbf{Phần 1. Trắc nghiệm lựa chọn}
\setcounterpageref{ex}{0}

\begin{ex} %[3D2V12]
Tính giá trị của biểu thức: $D = 9 \times 8 - 63 : 9 + 17$.
\choice
{$78$}
{$84$}
{$80$}
{\True $82$}
\loigiai{
Thực hiện nhân chia trước: $D = 72 - 7 + 17 = 65 + 17 = 82$.
\textbf{Đáp án của câu hỏi là D}
}
\end{ex}

\begin{ex} %[3D2V12]
Một lớp học có $36$ học sinh chia đều thành 9 nhóm. Hỏi 5 nhóm học tập đó có tất cả bao nhiêu học sinh?
\choice
{$18$ học sinh}
{\True $20$ học sinh}
{$24$ học sinh}
{$16$ học sinh}
\loigiai{
Mỗi nhóm có: $36 : 9 = 4$ học sinh.\
5 nhóm có: $5 \times 4 = 20$ học sinh.
\textbf{Đáp án của câu hỏi là B}
}
\end{ex}

\begin{ex} %[3D2V12]
Tìm $x$ biết: $x : 9 + 28 = 35$.
\choice
{$54$}
{$72$}
{\True $63$}
{$45$}
\loigiai{
Ta có: $x : 9 = 35 - 28 = 7 \Rightarrow x = 7 \times 9 = 63$.
\textbf{Đáp án của câu hỏi là C}
}
\end{ex}

\begin{ex} %[3D2V12]
Trong một phép chia cho $9$, thương bằng $8$ và số dư là số dư lớn nhất có thể có. Số bị chia là
\choice
{\True $80$}
{$79$}
{$81$}
{$78$}
\loigiai{
Vì số chia là 9 nên số dư lớn nhất là 8.\
Số bị chia là: $9 \times 8 + 8 = 72 + 8 = 80$.
\textbf{Đáp án của câu hỏi là A}
}
\end{ex}

\begin{ex} %[3D2V12]
Tính nhanh giá trị biểu thức: $9 \times 3 + 9 \times 4 + 9 \times 3$.
\choice
{$81$}
{$99$}
{$85$}
{\True $90$}
\loigiai{
Đặt thừa số chung: $9 \times (3 + 4 + 3) = 9 \times 10 = 90$.
\textbf{Đáp án của câu hỏi là D}
}
\end{ex}

\begin{ex} %[3D2V12]
Bác Hòa có $72\text{ kg}$ bưởi chia đều vào 9 thùng. Bác đã bán đi 4 thùng bưởi. Hỏi bác Hòa còn lại bao nhiêu ki-lô-gam bưởi?
\choice
{$36\text{ kg}$}
{\True $40\text{ kg}$}
{$45\text{ kg}$}
{$32\text{ kg}$}
\loigiai{
Mỗi thùng có: $72 : 9 = 8\text{ kg}$.\
Số thùng còn lại: $9 - 4 = 5$ thùng.\
Số bưởi còn lại: $5 \times 8 = 40\text{ kg}$.
\textbf{Đáp án của câu hỏi là B}
}
\end{ex}

\begin{ex} %[3D2V12]
Cho dãy số: $9; 18; 27; 36; \dots$ Số hạng thứ 9 của dãy số trên là
\choice
{$72$}
{$90$}
{\True $81$}
{$78$}
\loigiai{
Số hạng thứ 9 là: $9 \times 9 = 81$.
\textbf{Đáp án của câu hỏi là C}
}
\end{ex}

\begin{ex} %[3D2V12]
Tìm $x$ biết: $(x - 18) \times 9 = 54$.
\choice
{\True $24$}
{$22$}
{$26$}
{$25$}
\loigiai{
Ta có: $x - 18 = 54 : 9 = 6 \Rightarrow x = 6 + 18 = 24$.
\textbf{Đáp án của câu hỏi là A}
}
\end{ex}

\begin{ex} %[3D2V12]
Một đàn gà có $9$ con gà trống. Số con gà mái gấp $5$ lần số con gà trống. Hỏi đàn gà có tất cả bao nhiêu con gà?
\choice
{$45$ con gà}
{$50$ con gà}
{$60$ con gà}
{\True $54$ con gà}
\loigiai{
Số gà mái là: $9 \times 5 = 45$ con.\
Tổng số gà là: $9 + 45 = 54$ con.
\textbf{Đáp án của câu hỏi là D}
}
\end{ex}

\begin{ex} %[3D2V12]
Điền dấu thích hợp: $9 \times 6 - 15 \dots 81 : 9 \times 4$.
\choice
{$<$}
{\True $>$}
{$=$}
{Không so sánh được}
\loigiai{
Vế trái: $9 \times 6 - 15 = 54 - 15 = 39$.\
Vế phải: $81 : 9 \times 4 = 9 \times 4 = 36$.\
Vì $39 > 36$ nên điền dấu $>$. 
\textbf{Đáp án của câu hỏi là B}
}
\end{ex}

\begin{ex} %[3D2V12]
Một thanh sắt dài $81\text{ dm}$ được cắt thành các đoạn dài $9\text{ dm}$. Hỏi người thợ phải cưa tất cả bao nhiêu lần?
\choice
{$9$ lần cưa}
{$7$ lần cưa}
{\True $8$ lần cưa}
{$10$ lần cưa}
\loigiai{
Số đoạn sắt là: $81 : 9 = 9$ (đoạn).\
Số lần cưa cần thực hiện: $9 - 1 = 8$ (lần).
\textbf{Đáp án của câu hỏi là C}
}
\end{ex}

\begin{ex} %[3D2V12]
Tìm số tự nhiên $a$ biết: $9 \times 7 < a < 9 \times 9$ và $a$ là số tròn chục.
\choice
{\True $70$}
{$80$}
{$60$}
{$90$}
\loigiai{
Ta có: $63 < a < 81$. Các số tròn chục thỏa mãn là $70$ và $80$. Nếu là số tròn chục duy nhất: đổi đề hoặc chọn 70 (khi $a$ có hàng đơn vị là 0 và hàng chục là 7 $\Rightarrow a = 70$).
\textbf{Đáp án của câu hỏi là A}
}
\end{ex}

\textbf{Phần 2. Trắc nghiệm đúng sai}
\setcounterpageref{ex}{0}

\begin{ex} %[3D2V12]
Xét tính đúng sai của các nhận định sau về số học và phép nhân 9:
\choiceTF
{\True Tổng các chữ số của tích bất kì trong bảng nhân 9 luôn bằng 9 hoặc chia hết cho 9.}
{Mọi số chia hết cho 9 đều là số lẻ.}
{\True Tích $9 \times 6$ bằng tích $6 \times 9$.}
{\True Trong phép chia cho 9, số dư lớn nhất có thể là 8.}
\loigiai{
\begin{itemize}
    \item Ý a) Đúng: Ví dụ: 18 (1+8=9), 27 (2+7=9), 81 (8+1=9), 90 (9+0=9).
    \item Ý b) Sai: Ví dụ 18, 36, 54 chia hết cho 9 nhưng là số chẵn.
    \item Ý c) Đúng: Tính chất giao hoán của phép nhân.
    \item Ý d) Đúng: Số dư luôn bé hơn số chia.
\end{itemize}
\begin{itemize}
    \item a) \textbf{Đúng}.
    \item b) \textbf{Sai}.
    \item c) \textbf{Đúng}.
    \item d) \textbf{Đúng}.
\end{itemize}
}
\end{ex}

\begin{ex} %[3D2V12]
Một người mua 9 gói kẹo, mỗi gói có 8 cái kẹo. Người đó đã cho đi 24 cái kẹo:
\choiceTF
{\True Ban đầu người đó có tất cả 72 cái kẹo.}
{\True Sau khi cho đi, người đó còn lại 48 cái kẹo.}
{Số kẹo còn lại đủ để chia đều cho 5 bạn mỗi bạn 10 cái.}
{\True Số kẹo đã cho đi tương ứng với 3 gói kẹo ban đầu.}
\loigiai{
\begin{itemize}
    \item Ý a) Đúng: $9 \times 8 = 72$ cái kẹo.
    \item Ý b) Đúng: $72 - 24 = 48$ cái kẹo.
    \item Ý c) Sai: Cần $5 \times 10 = 50$ cái, mà chỉ còn 48 cái.
    \item Ý d) Đúng: $3 \times 8 = 24$ cái kẹo.
\end{itemize}
\begin{itemize}
    \item a) \textbf{Đúng}.
    \item b) \textbf{Đúng}.
    \item c) \textbf{Sai}.
    \item d) \textbf{Đúng}.
\end{itemize}
}
\end{ex}

\textbf{Phần 3. Trả lời ngắn}
\setcounterpageref{ex}{0}

\begin{ex} %[3D2V12]
Tính giá trị của biểu thức: $9 \times 7 - 45 : 9$.
\shortans{58}
\loigiai{
\textbf{Đáp số:} $58$\\
$9 \times 7 - 45 : 9 = 63 - 5 = 58$.
}
\end{ex}

\begin{ex} %[3D2V12]
Tìm số bị chia trong phép chia cho 9, biết thương là 7 và số dư là 6.
\shortans{69}
\loigiai{
\textbf{Đáp số:} $69$\\
$9 \times 7 + 6 = 63 + 6 = 69$.
}
\end{ex}

\begin{ex} %[3D2V12]
Một người có 63 quả cam chia đều vào 9 túi. Hỏi 6 túi như thế có bao nhiêu quả cam?
\shortans{42}
\loigiai{
\textbf{Đáp số:} $42$\\
Mỗi túi có $63 : 9 = 7$ quả. 6 túi có $6 \times 7 = 42$ quả.
}
\end{ex}

\begin{ex} %[3D2V12]
Tìm $x$ biết: $x : 9 = 16 - 8$.
\shortans{72}
\loigiai{
\textbf{Đáp số:} $72$\\
$x : 9 = 8 \Rightarrow x = 8 \times 9 = 72$.
}
\end{ex}

\begin{ex} %[3D2V12]
Một hình vuông có cạnh dài 9 cm. Tính chu vi hình vuông đó.
\shortans{36}
\loigiai{
\textbf{Đáp số:} $36$\\
Chu vi hình vuông là: $9 \times 4 = 36\text{ cm}$.
}
\end{ex}

\begin{ex} %[3D2V12]
Tính nhanh: $9 \times 1 + 9 \times 2 + 9 \times 3 + 9 \times 4$.
\shortans{90}
\loigiai{
\textbf{Đáp số:} $90$\\
$9 \times (1 + 2 + 3 + 4) = 9 \times 10 = 90$.
}
\end{ex}
